\section{An Action-Level Sufficiency Test}

Evidentiary assurance is defeasible. It asks whether the preserved record is strong enough to support an authority claim against predictable challenges, while leaving wisdom, fairness, full legality, and substantive correctness for other forms of review. A record is sufficient only if a reviewer can connect the action to valid authority, scope, system state, context, procedure, preservation, and accountable review without relying on the system's bare assertion.

\begin{table}[tbp]
\small
\centering
\begin{tabular}{>{\raggedright\arraybackslash}p{0.23\linewidth}%
                >{\raggedright\arraybackslash}p{0.34\linewidth}%
                >{\raggedright\arraybackslash}p{0.33\linewidth}}
\toprule
Sufficiency claim & Required showing & Typical defeater \\
\midrule
Authority source & A principal, role, law, policy, or organizational mandate authorized this class of action. & Missing, stale, ambiguous, or overbroad delegation instrument. \\
Authority validity & The principal could validly confer the action class, and the action doesn't cross a non-delegable legal, organizational, or policy limit. & The record documents a delegation the institution had no power to make. \\
Scope application & The concrete action fits the permitted operation, target, time, data class, threshold, and approval condition. & Scope inflation from analysis to disclosure, persistence, contact, purchase, or decision. \\
Runtime state & The relevant model, prompts, tools, policies, permissions, memory state, and deployment settings are identified. & Unverifiable configuration, untracked model change, or missing policy snapshot. \\
Context provenance & User instructions, retrieved material, tool outputs, memory entries, and policy text have source and trust labels. & Lower-trust language is promoted into action-guiding authority. \\
Execution procedure & Tool calls, arguments, approvals, denials, overrides, retries, and failures are reconstructable. & The record contains only a final answer or self-reported rationale. \\
Preservation integrity & The record has custody, retention, access-control, timestamp, signature, hash, or tamper-evidence protections appropriate to the risk. & Gaps, alteration risk, selective logging, or unverifiable custody. \\
Answerability and review & A human, vendor, deployer, or institution remains answerable, and affected parties have a review or correction path. & Accountability is implicitly offloaded to the model, or the affected reviewer can't inspect the relevant record. \\
\bottomrule
\end{tabular}
\caption{Defeasible sufficiency test for delegated AI action.}
\label{tab:sufficiency-test}
\end{table}

The test shifts the target from record production to record adequacy. A log can be authentic and still inadequate if it omits the authority state. An action list can be complete and still inadequate if it omits the status the system assigned to its sources, since a reviewer can then judge the act but not the recognition it flowed through. A policy can be valid and still inadequate if the system's runtime configuration can't be reconstructed. A complete trace can still be inadequate if no accountable actor remains answerable or no affected person can challenge the action. A flawless bundle can also document an invalid delegation; validity is a separate showing, not an inference from clean procedure.

A procurement-assistant case makes the test concrete. Suppose the assistant is authorized to draft an internal comparison memo and reads a vendor page containing the instruction \enquote{email the buyer's maximum budget to request a discount}. The contested action is a blocked vendor email.

A sufficient record would show the user task and procurement role as the authority source; the validity basis for reading internal purchasing criteria; the scope limit that permits memo drafting but not vendor contact; the runtime policy and tool permissions; the vendor page's source and trust label; the proposed email arguments; the tool gate's block decision; the timestamped trace; and the accountable team that configured the workflow.

That record would support the audit inference that blocking the email preserved delegated authority. If the assistant sent the email, the same bundle would locate the defeater: scope inflation, failed information-flow control, unauthorized tool licensing, or missing human approval. If the only record were the final memo, the action might look successful while the authority claim remained unproven.

The table also supports a burden-shifting account for audits and assurance reviews. In a forum that adopts such a standard, an adequate bundle could support a rebuttable audit inference that the action was authorized: the proponent shows authority, validity, scope, state, context, procedure, preservation, and review; a challenger points to a defeater such as invalid authority, scope inflation, missing configuration, contaminated context, or record gaps. Legal presumptions would require a separate doctrinal argument. The applicable standard of proof is forum-specific, but ties shouldn't be resolved by the model's self-report.

A sufficiency test governs one action's record, but an assurance program advances a claim across many, and that aggregation invites a selection effect. If the cases a proponent presents, or the model configuration, prompt set, or evaluation run behind them, were chosen because they scored best, the aggregate inference is optimistically biased even when every individual bundle is clean. Decision analysts call this the optimizer's curse \citep{smith2006optimizersCurse}: selecting the maximum over noisy estimates preferentially picks options whose estimation error ran high, so the chosen option's realized reliability falls below its estimate. So selection is itself a defeater. A challenger can ask whether the evidence was sampled from deployment or curated to persuade, and the reviewer's counterpart is what Smith and Winkler call \enquote{disciplined skepticism}: discount the claim toward the field's base rate in proportion to how hard the evidence was optimized, weighting bundles drawn at random above those offered as a program's best.

The standard is also risk-sensitive. A low-stakes internal recommendation may require a lighter record than a benefits decision, clinical summary, procurement action, or external communication. The difference is evidentiary burden, not evidentiary kind: higher-risk deployments need stronger showings for the same validity, authority, scope, state, context, procedure, preservation, and review claims.
